\documentclass{article}
\usepackage{fullpage}

% ===== BASIC PACKAGES =====
\usepackage{amsmath}
\usepackage{amsthm}
\usepackage{amsfonts}
\usepackage{amssymb}
\usepackage{mathtools}
\usepackage{microtype}
\usepackage{xfrac}
\usepackage{bbm}

\newcommand\numberthis{\addtocounter{equation}{1}\tag{\theequation}}

% ===== GRAPHICS AND FIGURES =====
\usepackage{graphicx}
\usepackage{subcaption}
\usepackage{float}
\usepackage{wrapfig}

% ===== TABLES =====
\usepackage{booktabs}
\usepackage{multirow}
\usepackage{makecell}
\usepackage{array}

% ===== ALGORITHMS =====
\usepackage{algorithm}
\usepackage{algorithmic}
% Alternative: \usepackage{algpseudocode} 

% ===== COLORS AND HYPERLINKS =====
\usepackage{xcolor}
\usepackage{hyperref}
\hypersetup{
    colorlinks=true,
    linkcolor=blue,
    citecolor=blue,
    urlcolor=blue
}

% ===== BIBLIOGRAPHY =====
\PassOptionsToPackage{numbers, compress}{natbib}


% ===== ADDITIONAL USEFUL PACKAGES =====
\usepackage{enumitem}  % Better control over lists
\usepackage{bbm}       % For indicator functions \mathbbm{1}
\usepackage{dsfont}    % Alternative math fonts
\usepackage{nicefrac}  % Nice inline fractions
\usepackage{bm}        % Bold math symbols
\usepackage{centernot} % For negated symbols


% ========================================
% MATHEMATICAL OPERATORS
% ========================================

% ===== LINEAR ALGEBRA =====
\DeclareMathOperator{\Tr}{Tr}
\DeclareMathOperator{\rank}{rank}
\DeclareMathOperator{\diag}{diag}
\DeclareMathOperator{\Diag}{Diag}
\DeclareMathOperator{\vect}{vec}
\DeclareMathOperator{\supp}{supp}
\DeclareMathOperator{\spn}{span}
\DeclareMathOperator{\Vol}{Vol}

% ===== OPTIMIZATION =====
\DeclareMathOperator*{\argmax}{arg\,max}
\DeclareMathOperator*{\argmin}{arg\,min}
\DeclareMathOperator*{\minimize}{minimize}
\DeclareMathOperator*{\maximize}{maximize}
\DeclareMathOperator{\prox}{prox}
\DeclareMathOperator{\grad}{grad}
\DeclareMathOperator{\Hess}{Hess}

% ===== PROBABILITY AND STATISTICS =====
\DeclareMathOperator{\Var}{Var}
\DeclareMathOperator{\Cov}{Cov}
\DeclareMathOperator{\Corr}{Corr}
\DeclareMathOperator{\sign}{sign}
\DeclareMathOperator*{\esssup}{ess\,sup}
\DeclareMathOperator*{\essinf}{ess\,inf}

% ===== NEURAL NETWORK SPECIFIC =====
\DeclareMathOperator{\ReLU}{ReLU}
\DeclareMathOperator{\softmax}{softmax}
\DeclareMathOperator{\sigmoid}{sigmoid}

\DeclareMathOperator{\width}{width}
\DeclareMathOperator{\depth}{depth}


% ========================================
% NOTATION: VECTORS AND MATRICES
% ========================================

% ===== LOWERCASE BOLD (vectors) =====
\newcommand{\ba}{\mathbf{a}}
\newcommand{\bb}{\mathbf{b}}
\newcommand{\bc}{\mathbf{c}}
\newcommand{\bd}{\mathbf{d}}
\newcommand{\be}{\mathbf{e}}
\newcommand{\bff}{\mathbf{f}}
\newcommand{\bg}{\mathbf{g}}
\newcommand{\bh}{\mathbf{h}}
\newcommand{\boldm}{\mathbf{m}}
\newcommand{\bp}{\mathbf{p}}
\newcommand{\bq}{\mathbf{q}}
\newcommand{\br}{\mathbf{r}}
\newcommand{\bs}{\mathbf{s}}
\newcommand{\bt}{\mathbf{t}}
\newcommand{\bu}{\mathbf{u}}
\newcommand{\bv}{\mathbf{v}}
\newcommand{\bw}{\mathbf{w}}
\newcommand{\bx}{\mathbf{x}}
\newcommand{\by}{\mathbf{y}}
\newcommand{\bz}{\mathbf{z}}

% ===== UPPERCASE BOLD (matrices) =====
\newcommand{\bA}{\mathbf{A}}
\newcommand{\bB}{\mathbf{B}}
\newcommand{\bC}{\mathbf{C}}
\newcommand{\bD}{\mathbf{D}}
\newcommand{\bE}{\mathbf{E}}
\newcommand{\bF}{\mathbf{F}}
\newcommand{\bG}{\mathbf{G}}
\newcommand{\bH}{\mathbf{H}}
\newcommand{\bI}{\mathbf{I}}
\newcommand{\bJ}{\mathbf{J}}
\newcommand{\bK}{\mathbf{K}}
\newcommand{\bL}{\mathbf{L}}
\newcommand{\bM}{\mathbf{M}}
\newcommand{\bN}{\mathbf{N}}
\newcommand{\bO}{\mathbf{O}}
\newcommand{\bP}{\mathbf{P}}
\newcommand{\bQ}{\mathbf{Q}}
\newcommand{\bR}{\mathbf{R}}
\newcommand{\bS}{\mathbf{S}}
\newcommand{\bT}{\mathbf{T}}
\newcommand{\bU}{\mathbf{U}}
\newcommand{\bV}{\mathbf{V}}
\newcommand{\bW}{\mathbf{W}}
\newcommand{\bX}{\mathbf{X}}
\newcommand{\bY}{\mathbf{Y}}
\newcommand{\bZ}{\mathbf{Z}}

% ===== SPECIAL VECTORS =====
\newcommand{\bzero}{\mathbf{0}}
\newcommand{\ones}{\mathbf{1}}

% ===== GREEK BOLD =====
\newcommand{\balpha}{\boldsymbol{\alpha}}
\newcommand{\bbeta}{\boldsymbol{\beta}}
\newcommand{\bgamma}{\boldsymbol{\gamma}}
\newcommand{\bdelta}{\boldsymbol{\delta}}
\newcommand{\bepsilon}{\boldsymbol{\epsilon}}
\newcommand{\bvarepsilon}{\boldsymbol{\varepsilon}}
\newcommand{\btheta}{\boldsymbol{\theta}}
\newcommand{\blambda}{\boldsymbol{\lambda}}
\newcommand{\bmu}{\boldsymbol{\mu}}
\newcommand{\bnu}{\boldsymbol{\nu}}
\newcommand{\bxi}{\boldsymbol{\xi}}
\newcommand{\bpi}{\boldsymbol{\pi}}
\newcommand{\brho}{\boldsymbol{\rho}}
\newcommand{\bsigma}{\boldsymbol{\sigma}}
\newcommand{\btau}{\boldsymbol{\tau}}
\newcommand{\bphi}{\boldsymbol{\phi}}
\newcommand{\bvarphi}{\boldsymbol{\varphi}}
\newcommand{\bpsi}{\boldsymbol{\psi}}
\newcommand{\bomega}{\boldsymbol{\omega}}

\newcommand{\bGamma}{\boldsymbol{\Gamma}}
\newcommand{\bDelta}{\boldsymbol{\Delta}}
\newcommand{\bTheta}{\boldsymbol{\Theta}}
\newcommand{\bLambda}{\boldsymbol{\Lambda}}
\newcommand{\bSigma}{\boldsymbol{\Sigma}}
\newcommand{\bUpsilon}{\boldsymbol{\Upsilon}}
\newcommand{\bPhi}{\boldsymbol{\Phi}}
\newcommand{\bPsi}{\boldsymbol{\Psi}}
\newcommand{\bOmega}{\boldsymbol{\Omega}}
\newcommand{\bPi}{\boldsymbol{\Pi}}

% ===== TILDE NOTATION =====
\newcommand{\tA}{\widetilde{A}}
\newcommand{\tB}{\widetilde{B}}
\newcommand{\tL}{\widetilde{L}}
\newcommand{\tU}{\widetilde{U}}
\newcommand{\tV}{\widetilde{V}}
\newcommand{\tS}{\widetilde{\Sigma}}
\newcommand{\ts}{\widetilde{\sigma}}
\newcommand{\tW}{\widetilde{W}}
\newcommand{\tX}{\widetilde{X}}
\newcommand{\tY}{\widetilde{Y}}
\newcommand{\tf}{\widetilde{f}}
\newcommand{\tg}{\widetilde{g}}


% ========================================
% NORMS AND METRICS
% ========================================
\newcommand{\norm}[1]{\left\lVert#1\right\rVert}
\newcommand{\abs}[1]{\left\lvert#1\right\rvert}
\newcommand{\normone}[1]{\left\lVert#1\right\rVert_1}
\newcommand{\normtwo}[1]{\left\lVert#1\right\rVert_2}
\newcommand{\normF}[1]{\left\lVert#1\right\rVert_\text{F}}
\newcommand{\norminf}[1]{\left\lVert#1\right\rVert_\infty}
\newcommand{\normnuc}[1]{\left\lVert#1\right\rVert_*}
\newcommand{\normop}[1]{\left\lVert#1\right\rVert_\text{op}}
\newcommand{\normspec}[1]{\left\lVert#1\right\rVert_2}
\newcommand{\inner}[2]{\left\langle#1,#2\right\rangle}


% ========================================
% PROBABILITY NOTATION
% ========================================
\renewcommand{\Pr}{\mathbb{P}}
\newcommand{\E}{\mathbb{E}}

\newcommand{\given}{\,|\,}
\newcommand{\ind}{\mathbbm{1}}  % Indicator function
\newcommand{\iid}{\text{i.i.d.}}
\newcommand{\as}{\text{a.s.}}


% ========================================
% SETS AND SPACES
% ========================================
\newcommand{\R}{\mathbb{R}}
\newcommand{\C}{\mathbb{C}}
\newcommand{\N}{\mathbb{N}}
\newcommand{\Z}{\mathbb{Z}}
\newcommand{\Q}{\mathbb{Q}}
\newcommand{\cA}{\mathcal{A}}
\newcommand{\cB}{\mathcal{B}}
\newcommand{\cC}{\mathcal{C}}
\newcommand{\cD}{\mathcal{D}}
\newcommand{\cE}{\mathcal{E}}
\newcommand{\cF}{\mathcal{F}}
\newcommand{\cG}{\mathcal{G}}
\newcommand{\cH}{\mathcal{H}}
\newcommand{\cL}{\mathcal{L}}
\newcommand{\cM}{\mathcal{M}}
\newcommand{\cN}{\mathcal{N}}
\newcommand{\cO}{\mathcal{O}}
\newcommand{\cP}{\mathcal{P}}
\newcommand{\cS}{\mathcal{S}}
\newcommand{\cT}{\mathcal{T}}
\newcommand{\cU}{\mathcal{U}}
\newcommand{\cV}{\mathcal{V}}
\newcommand{\cW}{\mathcal{W}}
\newcommand{\cX}{\mathcal{X}}
\newcommand{\cY}{\mathcal{Y}}
\newcommand{\cZ}{\mathcal{Z}}


% ========================================
% NEURAL NETWORK SPECIFIC NOTATION
% ========================================

% ===== NETWORK PARAMETERS =====
\newcommand{\layer}[1]{^{[#1]}}           % Layer superscript
\newcommand{\Wlayer}[1]{\bW\layer{#1}}    % Weight matrix for layer
\newcommand{\blayer}[1]{\mathbf{b}\layer{#1}}  % Bias vector for layer
\newcommand{\alayer}[1]{\mathbf{a}\layer{#1}}  % Activation for layer
\newcommand{\zlayer}[1]{\mathbf{z}\layer{#1}}  % Pre-activation for layer

% ===== COMMON NEURAL NETWORK NOTATION =====
\newcommand{\loss}{\ell}                   % Loss function
\newcommand{\risk}{\mathcal{R}}           % Risk/Expected loss
\newcommand{\emprisk}{\widehat{\mathcal{R}}}  % Empirical risk
\newcommand{\model}{\cH}                  % Hypothesis class/model class


% ========================================
% GENERALIZATION BOUNDS NOTATION
% ========================================

% ===== COMPLEXITY MEASURES =====
\newcommand{\Rad}{\mathfrak{R}}           % Rademacher complexity
\newcommand{\EmpRad}{\widehat{\mathfrak{R}}}  % Empirical Rademacher complexity
\newcommand{\VC}{\text{VC}}               % VC dimension
\newcommand{\cover}{\mathcal{N}}          % Covering number
\newcommand{\pack}{\mathcal{M}}           % Packing number
\newcommand{\fat}{\text{fat}}             % Fat-shattering dimension

% ===== COMPRESSION AND SPARSITY =====
\newcommand{\sparsity}{\text{sparsity}}
\newcommand{\compression}{\text{CR}}     % Compression ratio
\newcommand{\nnz}{\text{nnz}}            % Number of non-zeros


% ========================================
% OPTIMIZATION NOTATION
% ========================================
\newcommand{\learning}{\eta}              % Learning rate
\newcommand{\momentum}{\beta}             % Momentum parameter
\newcommand{\lip}{L}                      % Lipschitz constant
\newcommand{\smooth}{\beta}               % Smoothness constant
\newcommand{\strongly}{\mu}               % Strong convexity parameter
\newcommand{\hess}{\nabla^2}              % Hessian


% ========================================
% COMPLEXITY NOTATION
% ========================================
\newcommand{\BigO}{\mathcal{O}}
\newcommand{\BigOmega}{\Omega}
\newcommand{\BigTheta}{\Theta}
\newcommand{\softO}{\widetilde{\mathcal{O}}}
\newcommand{\poly}{\text{poly}}


% ========================================
% SPECIAL NOTATION
% ========================================
\newcommand{\defeq}{\coloneqq}            % Definition equals
\newcommand{\eqdef}{\eqqcolon}            % Equals by definition



% ========================================
% THEOREM ENVIRONMENTS
% ========================================
\newtheorem{theorem}{Theorem}[section]
\newtheorem{lemma}[theorem]{Lemma}
\newtheorem{proposition}[theorem]{Proposition}
\newtheorem{corollary}[theorem]{Corollary}
\newtheorem{claim}[theorem]{Claim}

\theoremstyle{definition}
\newtheorem{definition}[theorem]{Definition}
\newtheorem{example}[theorem]{Example}
\newtheorem{assumption}[theorem]{Assumption}

\theoremstyle{remark}
\newtheorem{remark}[theorem]{Remark}
\newtheorem{note}[theorem]{Note}
\newtheorem{observation}[theorem]{Observation}


% ========================================
% AUTHOR COMMENTS
% ========================================
\newcommand{\rk}[1]{{\color{purple}{[RK: #1]}}}
\newcommand{\pd}[1]{{\color{blue}{[PD: #1]}}}
\newcommand{\kaiwen}[1]{{\color{red}{[KW: #1]}}}
\newcommand{\todo}[1]{{\color{orange}{[TODO: #1]}}}


% ========================================
% PROJECT-SPECIFIC COMMANDS
% ========================================

% Add your project-specific notation here
\newcommand{\Reff}{\ensuremath{R_\text{eff}}}
\newcommand{\punif}{\ensuremath{p^\text{unif}}}


% ========================================
% MISCELLANEOUS
% ========================================

% Better spacing in align environments
\setlength{\jot}{8pt}

% Theorem-like environment without numbering for remarks
\newtheorem*{remark*}{Remark}
\newtheorem*{note*}{Note}